% GENERATED FILE. Edit canonical lessons, then run npm run generate:books.
% canonical-source-hash: fnv1a64:1c17af9ccd6e38e5
% canonical-lessons: LA-C525-tussio, LA-C525-vomo, LA-C525-sitio, LA-C525-lippus, LA-C525-lippitudo

\chapter{Coughs and sore eyes}
\label{ch:la-coughs-and-sore-eyes}
\hlchaptermodality{525}

\begin{chapteropening}
\textbf{By the end of this chapter:} \emph{I can describe coughing, sickness, thirst and sore eyes.}

Complete the last lesson of chapter 525: I can describe coughing, sickness, thirst and sore eyes.
\end{chapteropening}

\section[tussiō, tussīre]{tussiō, tussīre — to cough}

\label{lesson:LA-C525-tussio}



\begin{quote}
Before the new one: say the Latin for a sore, then the Latin for to swell.
\end{quote}

\subsection*{You'll want to know: tussiō, tussīre}
\textbf{tussiō, tussīre} — \textquotedblleft{}to cough\textquotedblright{}.

\begin{quote}
Senex tōtam noctem tussit.
\end{quote}

\begin{quote}
\emph{The old man coughs all night.}
\end{quote}

\emph{tussiō, tussīre}, from \emph{tussis} — a cough, already known.

\begin{cousinweb}
\textbf{tussiō} is built on \emph{tussis} \textquotedblleft{}a cough\textquotedblright{} (Lewis \& Short; EtymDB), which EtymDB traces to a Proto-Indo-European root for coughing. English \emph{pertussis}, whooping cough, comes from the noun.
\end{cousinweb}

\begin{grammarlens}[title={What you've built}]
One more word for telling, explaining and describing at B1.
\end{grammarlens}

\subsection*{Your turn}
\begin{itemize}
\raggedright
  \item \emph{Say it:} \emph{tussiō, tussīre}
  \item \emph{Say it:} \emph{tussiō, tussīre}, once more
  \item \emph{Say it:} a sentence of your own with \emph{tussiō, tussīre}
  \item \emph{Recall:} say the Latin for a sore, then the Latin for to swell, then say \emph{tussiō, tussīre} again
\end{itemize}

\begin{tcolorbox}[breakable,colback=teal!4,colframe=teal!35!black,title={Before you move on}]
What does tussiō, tussīre mean? (\textquotedblleft{}To cough\textquotedblright{}.) Say it once more.
\end{tcolorbox}
\section[vomō, vomere]{vomō, vomere — to vomit, be sick}

\label{lesson:LA-C525-vomo}



\begin{quote}
Before the new one: say the Latin for to swell, then the Latin for to cough.
\end{quote}

\subsection*{You'll want to know: vomō, vomere}
\textbf{vomō, vomere} — \textquotedblleft{}to vomit, be sick\textquotedblright{}.

\begin{quote}
Nauta in nāve vomit.
\end{quote}

\begin{quote}
\emph{The sailor is sick on the ship.}
\end{quote}

\emph{vomō, vomere, vomuī, vomitum}. Also of a volcano spewing fire. English \emph{vomit}.

\begin{cousinweb}
\textbf{vomō} is inherited from a Proto-Indo-European root meaning \textquotedblleft{}to vomit\textquotedblright{}, and Greek \emph{emeo} is a cousin (EtymDB; Lewis \& Short). English \emph{vomit} comes from it.
\end{cousinweb}

\begin{grammarlens}[title={What you've built}]
One more word for telling, explaining and describing at B1.
\end{grammarlens}

\subsection*{Your turn}
\begin{itemize}
\raggedright
  \item \emph{Say it:} \emph{vomō, vomere}
  \item \emph{Say it:} \emph{vomō, vomere}, once more
  \item \emph{Say it:} a sentence of your own with \emph{vomō, vomere}
  \item \emph{Recall:} say the Latin for to swell, then the Latin for to cough, then say \emph{vomō, vomere} again
\end{itemize}

\begin{tcolorbox}[breakable,colback=teal!4,colframe=teal!35!black,title={Before you move on}]
What does vomō, vomere mean? (\textquotedblleft{}To vomit, be sick\textquotedblright{}.) Say it once more.
\end{tcolorbox}
\section[sitiō, sitīre]{sitiō, sitīre — to be thirsty}

\label{lesson:LA-C525-sitio}



\begin{quote}
Before the new one: say the Latin for to cough, then the Latin for to vomit.
\end{quote}

\subsection*{You'll want to know: sitiō, sitīre}
\textbf{sitiō, sitīre} — \textquotedblleft{}to be thirsty\textquotedblright{}.

\begin{quote}
Post cursum valdē sitiō.
\end{quote}

\begin{quote}
\emph{After running I am very thirsty.}
\end{quote}

\emph{sitiō, sitīre, sitīvī}, from \emph{sitis} — thirst, already known.

\begin{cousinweb}
\textbf{sitiō} is built on \emph{sitis} \textquotedblleft{}thirst\textquotedblright{} (Lewis \& Short; EtymDB).
\end{cousinweb}

\begin{grammarlens}[title={What you've built}]
One more word for telling, explaining and describing at B1.
\end{grammarlens}

\subsection*{Your turn}
\begin{itemize}
\raggedright
  \item \emph{Say it:} \emph{sitiō, sitīre}
  \item \emph{Say it:} \emph{sitiō, sitīre}, once more
  \item \emph{Say it:} a sentence of your own with \emph{sitiō, sitīre}
  \item \emph{Recall:} say the Latin for to cough, then the Latin for to vomit, then say \emph{sitiō, sitīre} again
\end{itemize}

\begin{tcolorbox}[breakable,colback=teal!4,colframe=teal!35!black,title={Before you move on}]
What does sitiō, sitīre mean? (\textquotedblleft{}To be thirsty\textquotedblright{}.) Say it once more.
\end{tcolorbox}
\section[lippus, lippa, lippum]{lippus, lippa, lippum — bleary-eyed, with sore eyes}

\label{lesson:LA-C525-lippus}



\begin{quote}
Before the new one: say the Latin for to vomit, then the Latin for to be thirsty.
\end{quote}

\subsection*{You'll want to know: lippus, lippa, lippum}
\textbf{lippus, lippa, lippum} — \textquotedblleft{}bleary-eyed, with sore eyes\textquotedblright{}.

\begin{quote}
Servus lippus librum legere nōn potest.
\end{quote}

\begin{quote}
\emph{The bleary-eyed slave cannot read the book.}
\end{quote}

Watery, inflamed eyes, a common complaint in Rome.

\begin{cousinweb}
\textbf{lippus} is a native Latin adjective; Lewis \& Short compares a root meaning \textquotedblleft{}to smear\textquotedblright{}, but the origin is not settled.
\end{cousinweb}

\begin{grammarlens}[title={What you've built}]
One more word for telling, explaining and describing at B1.
\end{grammarlens}

\subsection*{Your turn}
\begin{itemize}
\raggedright
  \item \emph{Say it:} \emph{lippus, lippa, lippum}
  \item \emph{Say it:} \emph{lippus, lippa, lippum}, once more
  \item \emph{Say it:} a sentence of your own with \emph{lippus, lippa, lippum}
  \item \emph{Recall:} say the Latin for to vomit, then the Latin for to be thirsty, then say \emph{lippus, lippa, lippum} again
\end{itemize}

\begin{tcolorbox}[breakable,colback=teal!4,colframe=teal!35!black,title={Before you move on}]
What does lippus, lippa, lippum mean? (\textquotedblleft{}Bleary-eyed, with sore eyes\textquotedblright{}.) Say it once more.
\end{tcolorbox}
\section[lippitūdō, lippitūdinis]{lippitūdō, lippitūdinis — inflammation of the eyes}

\label{lesson:LA-C525-lippitudo}



\begin{quote}
Before the new one: say the Latin for to be thirsty, then the Latin for bleary-eyed.
\end{quote}

\subsection*{You'll want to know: lippitūdō, lippitūdinis}
\textbf{lippitūdō, lippitūdinis} — \textquotedblleft{}inflammation of the eyes\textquotedblright{}.

\begin{quote}
Propter lippitūdinem epistulās nōn scrībō.
\end{quote}

\begin{quote}
\emph{Because of sore eyes I am not writing letters.}
\end{quote}

Feminine, third declension, from \emph{lippus}. A complaint of busy letter-writers.

\begin{cousinweb}
\textbf{lippitūdō} is built on \emph{lippus} \textquotedblleft{}bleary-eyed\textquotedblright{} with \emph{-tūdō} (Lewis \& Short).
\end{cousinweb}

\begin{grammarlens}[title={What you've built}]
One more word for telling, explaining and describing at B1.
\end{grammarlens}

\subsection*{Your turn}
\begin{itemize}
\raggedright
  \item \emph{Say it:} \emph{lippitūdō, lippitūdinis}
  \item \emph{Say it:} \emph{lippitūdō, lippitūdinis}, once more
  \item \emph{Say it:} a sentence of your own with \emph{lippitūdō, lippitūdinis}
  \item \emph{Recall:} say the Latin for to be thirsty, then the Latin for bleary-eyed, then say \emph{lippitūdō, lippitūdinis} again
\end{itemize}

\begin{tcolorbox}[breakable,colback=teal!4,colframe=teal!35!black,title={Before you move on}]
What does lippitūdō, lippitūdinis mean? (\textquotedblleft{}Inflammation of the eyes\textquotedblright{}.) Say it once more.
\end{tcolorbox}
